\documentclass[10pt,a4paper]{amsart}
\usepackage[margin=19mm]{geometry}
\usepackage{amsmath,amssymb,mathtools}
\usepackage[scr=rsfs,cal=euler]{mathalfa}
\usepackage{xfrac}
\usepackage[unicode,hidelinks,bookmarks=false]{hyperref}
\newcommand{\ie}{i.e.\ }
\newcommand{\cf}{cf.\ }
\newcommand{\resp}{respectively\ }
\newcommand{\Subsection}{\subsection}
\providecommand{\nobreakdash}{\nobreak\hbox{-}}
\setlength{\emergencystretch}{2em}
\multlinegap0pt
\usepackage{bm}
\usepackage{bbm}
\usepackage{dsfont}
\usepackage{xspace}
\usepackage{cancel}
\usepackage{mathtools}
\usepackage[shortlabels]{enumitem}
\usepackage{amsthm}
\makeatletter
\@ifundefined{defin}{\newtheorem{defin}{Defin}}{}
\@ifundefined{lemma}{\newtheorem{lemma}[defin]{Lemma}}{}
\@ifundefined{prop}{\newtheorem{prop}[defin]{Proposition}}{}
\@ifundefined{thm}{\newtheorem{thm}[defin]{Theorem}}{}
\makeatother
\newcommand{\Gammatilde}{\ensuremath{\widetilde{\Gamma}}}%
\newcommand{\acts}{\ensuremath{\curvearrowright}}%
\newcommand{\arr}{\longrightarrow}
\newcommand{\aut}{\ensuremath{\operatorname{Aut}}}%
\newcommand{\esf}{\ensuremath{\mathsf{e}}}
\newcommand{\rist}{\ensuremath{\operatorname{RiSt}}}%
\title[Commensurating actions and self-similar groups]{Commensurating actions and self-similar groups}
\author{Nicol\'as Matte Bon, Volodymyr Nekrashevych, Tianyi Zheng}
\date{Source: arXiv:2607.13776v2 (CC BY 4.0)}
\begin{document}
\maketitle

\noindent\emph{Source and licence.} Commensurating actions and self-similar groups. Version arXiv:2607.13776v2 (CC BY 4.0), distributed under the Creative Commons Attribution 4.0 International licence, as recorded in the arXiv licence metadata for this exact version and shown on the abstract page https://arxiv.org/abs/2607.13776v2. Full rights notice: licenses/arxiv-2607.13776-CC-BY-4.0.txt.\par\smallskip

\noindent\textit{Editorial scope.} Counted unit: the Thm whose source label is \texttt{t-multi-ended-graphs} in the article (source0.tex lines 530--537, source label \texttt{t-multi-ended-graphs}), delivered together with the article's own complete original proof of it (source0.tex lines 539--618, one \texttt{proof} environment of 80 lines). No mathematics is written, repaired or re-derived: statement and proof are the article's own text line for line. Required context retained uncounted: lem:endsofcoverings (lines 366--371); l-reduction-transitive (lines 380--385); t-commensurating-product-discrete (lines 395--401); pr:convergenceofgraphs (lines 463--466); l-quotient-rist (lines 509--513). Every definition, display and label of those lines is retained unchanged. Omitted from this package and not redistributed: 1--365, 372--379, 386--394, 402--462, 467--508, 514--529, 538--538, 619--1344. Nothing inside a retained excerpt is omitted or altered: every retained body is delivered exactly as the article's lines are written, so no omission receipt of any kind accompanies this package. Citations: the retained excerpts carry no citation command at all, so no bibliography entry of the article is transcribed and no reference is redistributed. Reference adaptation: none was needed and none was made; every reference inside the delivered unit resolves inside it, so no unresolved reference or citation is produced. Printed slips kept literal: no printed word, symbol, hypothesis or number of the counted statement or of its proof was altered. Layout and class: the article's own class is \texttt{amsart} with packages including amssymb, amsfonts, amsthm, amsmath, mathrsfs, xspace, hyperref, changes, todonotes, combelow, graphicx, babel, inputenc, fontenc; the wrapper is the lane's pinned \texttt{amsart} editorial wrapper at 10pt on A4, which restates the theorem-like environments the retained text needs and adds the article's own macro definitions that the retained text needs (\texttt{\textbackslash Gammatilde}, \texttt{\textbackslash acts}, \texttt{\textbackslash arr}, \texttt{\textbackslash aut}, \texttt{\textbackslash esf}, \texttt{\textbackslash rist}), with extra wrapper packages (bm, bbm, dsfont, xspace, cancel, mathtools, enumitem). The delivered numbering is the unit's own. Assets: no retained span contains a figure environment, a graphics-inclusion command, a PDF-page inclusion, a table or a TikZ picture; no image file is copied into this package, the package declares no asset dependency and no image file is redistributed with it. Licence: version arXiv:2607.13776v2 of this article is distributed under the Creative Commons Attribution 4.0 International licence, as recorded in the arXiv licence metadata for this exact version and shown on the abstract page https://arxiv.org/abs/2607.13776v2. A full rights notice is delivered at licenses/arxiv-2607.13776-CC-BY-4.0.txt. Characters: every retained excerpt is pure ASCII, so the delivered engine, XeLaTeX with its default Latin Modern text font, reports no missing character.
\begin{lemma}
\label{lem:endsofcoverings}
Let $G$ be a finitely generated group, and $K\le H\le G$ be subgroups such  that $[H:K]<\infty$. Then $\esf(\Gamma_{G/K})\ge \esf(\Gamma_{G/H})$. Furthermore, for every commensurated subset $A\subset G/H$, there exists a commensurated subset $\widetilde{A}\subset G/K$ such that $\nu_{\widetilde{A}}=[H: K] \nu_A$. 


\end{lemma}

\begin{lemma} \label{l-reduction-transitive}
	Let $G$ be a group and  $\alpha \colon G\acts X$  be transitive action having a commensurated subset $A\subset X$ which is not transfixed. Let $N$ be a finitely generated normal subgroup of $G$ such that $\alpha(N)$ is infinite. Then the following hold:
	\begin{enumerate}
	
	\item the action $\alpha \colon N \acts X$ has finitely many orbits;
	\item	 for every $N$-orbit $Y\subset X$ the sets $A\cap Y$ and $Y\setminus A$ are both infinite.\end{enumerate}	\end{lemma}

	\begin{thm} \label{t-commensurating-product-discrete}
	Let $G=G_1\times G_2$ where each $G_i$ is  finitely generated. Let $\alpha \colon G\acts X$ be a transitive action whose Schreier graph $\Gamma_X$ has at least two ends. Then  one of the following holds.
	\begin{enumerate}[label=(\roman*)]
	\item There exists $i=1, 2$ such that $\alpha(G_i)$ is finite. 
	\item   $\alpha(G)$ is virtually abelian.  
	\end{enumerate}
	\end{thm}

\begin{prop}
\label{pr:convergenceofgraphs}
 Let $\xi=(v_n)_n\in\partial T$. Then for every $R>0$ there exists $n$ such that the covering maps $\Gamma_\xi\arr\Gamma_n$ and $\Xi_n\arr\Gammatilde_\xi$ induce isomorphisms $B_{\Gamma_\xi}(\xi, R)\arr B_{\Gamma_n}(v_n, R)$ and $B_{\Xi_n}(G^0_{v_n}, R)\arr B_{\Gammatilde_\xi}([e,\xi], R)$.
\end{prop}

\begin{lemma} \label{l-quotient-rist}
    Let $G\le \aut(T)$ be a branch group, and fix $v\in T$. Then for every normal subgroup $N\unlhd \rist(v)$, there is $n\ge 1$ such that for every  $w\in T^{(n)}_v$, either $N$ contains $\rist(w)'$, or $N$ fixes $w$ and acts trivially on $T_w$.

    In particular $\rist(v)/N$ is virtually abelian if and only if $N$ contains $\rist(w)'$ for some $n\ge 1$ and all $w\in T^{(n)}_v$. 
\end{lemma}

\begin{thm} \label{t-multi-ended-graphs}
Let $G\le \aut(T)$ be a finitely generated branch group whose action on $\partial T$ has finite groups of germs.
Let $H\le G$ be a subgroup such that the action of $G$ on $G/H$ is faithful and $\esf(\Gamma_{G/H})\ge 2$. Then there is $\xi \in \partial T$ such that $H$ is a finite index subgroup of $G_\xi$. Furthermore, we have $\esf(\Gammatilde_\xi)\ge \esf(\Gamma_{G/H})$, and for every commensurated subset $A\subset \Gamma_{G/H}$, there is a commensurated subset $\widetilde{A}\subset \Gammatilde_
\xi$ such that $\nu_A-\alpha\nu_{\widetilde{A}}\in \ell^\infty(G)$, where $\alpha=\frac{[G^0_\xi\,  :\,  G^0_\xi\cap H]}{[H\, : \, G^0_\xi\cap H]}$.



\end{thm}

\begin{proof}
Fix a finite symmetric generating set $S$ of $G$, and let  $H\le G$ be a subgroup as in the statement. 


Given a geodesic ray $\xi\in \partial T$ in a rooted tree $T$, we write $N(\xi)\subset T$ for the collection of all vertices that are not on $\xi$, but whose parent is in $\xi$.   Note that we have	\[\partial T\setminus \{\xi\}=\bigsqcup_{v\in N(\xi)} \partial T_v.\] 
The starting point of the proof is the following application of Theorem \ref{t-commensurating-product-discrete} (which does not require the assumption that groups of germs are finite).

\begin{lemma}[First key step] \label{l-commensurating-branch}
There is $\xi\in \partial T$ and a finite index subgroup $K_v \le \rist(v)$ for each  $v\in  N(\xi)$ such that 
	\begin{equation}\label{e-subgroup} \bigoplus_{v\in N(\xi)} K_v \le H \le G_{\xi}.\end{equation}\end{lemma}
Here and throughout, we identify $\bigoplus_{v\in N(\xi)} K_v$ with the subgroup $\langle K_v, v\in N(\xi)\rangle$.
\begin{proof}[Proof of Lemma \ref{l-commensurating-branch}]
Let us set $X=G/H$, and let $A\subset X$ be a commensurated subset which is not transfixed. Fix $x\in X$, and let $n\ge 1$. Obviously the action of $\rist(n)$ on $X$ must have infinite image, since otherwise the action of $G$ would have finite image and thus $X$ would be finite. By Lemma \ref{l-reduction-transitive}, the action of $\rist(n)\simeq \prod_{v\in T^{(n)}}\rist(v)$ on $X$ has finitely many orbits, and for each such orbit $Y$ the sets $Y\cap A$ and $Y\cap A^c$ are both infinite.  Let $Y_n$ be the $\rist(n)$-orbit of $x$; it is infinite by Lemma \ref{l-reduction-transitive}. Hence there is at least one $v\in T^{(n)}$ such that $\rist(v)\acts Y_n$ has infinite image. Suppose, by contradiction, that there are two such distinct $v, w\in T^{(n)}$. Then we can write $\rist(n)=G_1\times G_2$, where $G_i\acts Y_n$ has infinite image for $i=1, 2$. Then  Theorem \ref{t-commensurating-product-discrete} implies that the action $\rist(n)\acts Y_n$ has virtually abelian image, and Lemma \ref{l-quotient-rist} implies that there is $m\ge n$ such that $\rist(m)'$ fixes every point of $Y_n$. But since $\rist(m)'$ is normal in $G$, and $G$ acts transitively on $X$, this implies that $\rist(m)'$ actually fixes every point of $X$, contradicting that $G$ acts faithfully on $X$.  Therefore, for each $n$, there exists a unique $v_n\in T^{(n)}$ such that the action of $\rist(v_n)$ on $Y_n$ has infinite image. The sequence $(v_n)$ forms a ray $\xi_x\in \partial T$. The map $x\mapsto \xi_x$ is equivariant, so that $G_x\le G_{\xi_x}$. In addition, for every $v\in N(\xi)$ at level $n$, we have $v\neq v_n$, so there is a finite index subgroup $K_v\le \rist(v)$ that acts trivially on $Y_{n}$, and thus fixes $x$. It follows that $\bigoplus_{v\in N(\xi_x)} K_v\le G_x\le G_{\xi_x}$. Choosing $x=eH$ we get the desired conclusion. \qedhere
\end{proof}



From now on, we let $\xi=(v_n)\in \partial T$ and $(K_v)_{v\in N(\xi)}$ be given by Lemma \ref{l-commensurating-branch}. We set $K=\bigoplus_{v\in N(\xi)} K_v$, so that \eqref{e-subgroup} reads $K\le H\le G_\xi$. In the reminder of the proof we assume that the group of germs $G_\xi /G^0_\xi$ is finite, and progressively strengthen the conclusion of Lemma \ref{l-commensurating-branch} until reaching the conclusion of the theorem. 


For each $v\in N(\xi)$, denote the set of maps $\{g|_{T_v}\colon T_v\to T_{g(v)} \,  \vert \, g\in G\}$ by $E_v$.
The group $K_v$ acts on $E_v$ from the right by precomposition, and the condition that $K_v$ has finite index in $\rist(v)$ implies that the quotient $E_v/K_v$ is a finite set. Then for every $g\in G$, the data 
\[(g|_{T_v}\text{ mod }K_v)_{v\in N(\xi)}
\in \prod_{v\in N(\xi)} E_v/K_v\]
determines uniquely the coset $g K$, and thus the coset $gH$. We will repeatedly use this observation. 
 
Since $G^0_\xi$ has finite index in $G_\xi\ge H$, the group $H\cap G^0_\xi$ has finite index in $H$. Moreover Lemma \ref{lem:endsofcoverings} ensures that $\esf(\Gamma_{G/(H\cap G^0_\xi)})\ge \esf(\Gamma_{G/H})$  and that for every commensurated subset $A\subset{G/H}$, there is a commensurated subset $B\subset {G/(H\cap G^0_\xi)}$  such that $\nu_A=[H\colon H\cap G^0_\xi]^{-1}\nu_B$.  Hence, to prove the theorem, there is no loss of generality in replacing $H$ by $H \cap G^0_\xi$, i.e. we can suppose that $H\le G^0_\xi$ to begin with. After this preliminary reduction,  the map $gH\mapsto gG^0_\xi$ induces a covering map 

\[p\colon \Gamma_{G/H}\to  \widetilde{\Gamma}_\xi.\]


\begin{lemma}\label{l-fiber-disconnected}
Every connected infinite subset $U\subset \Gamma_{G/H}$ has an infinite projection $p(U)$.
\end{lemma}
\begin{proof}
Suppose that $U$ is connected and $p(U)$ is finite, and let us show that $U$ is finite. Without loss of generality, we may assume that $U$ contains the basepoint $eH$ of $\Gamma_{G/H}$ (upon enlarging $U$ by adding to it a path from the basepoint to any point in $U$). Let $R>0$ be such that $p(U)\subset B_{\Gammatilde_\xi}(eG^0_\xi, R)$. By Proposition~\ref{pr:convergenceofgraphs}, for all $m$ large enough the map $g\cdot v\mapsto [g, \xi]$ is an isomorphism $B_{\Xi_m}(G^0_v, R)\arr B_{\Gammatilde_\xi}(G^0_\xi, R)$, where $v$ is the prefix of length $m$ of $\xi$.

Suppose that $s_1, \ldots, s_n\in S$ is any sequence of generators such that the path $\gamma=(eH, s_1H, s_2s_1H, \ldots, s_n\cdots s_1H)$ is contained in $U$, and thus its projection $p(\gamma)$ is contained in $B_{\Gammatilde_\xi}(G^0_\xi, R)$. 
Then the preimage of $p(\gamma)$ in $B_{\Xi_m}(eG^0_v, R)$ is the path 
$G^0_v, s_1G^0_v, s_2s_1G^0_v , \ldots, s_n\cdots s_1G^0_v$. Since it stays inside $B_{\Xi_m}(G^0_v, R)$, the maps $s_i\cdots s_2s_1|_{T_v}$ belong to a finite set.  The value of $s_n\cdots s_1\vert_{T_v}$ determines uniquely the map $s_n\cdots s_1|_{T_w}$ for all $w\in T_v$, and hence for all but finitely many $w\in N(\xi)$. We deduce that  $(s_n\cdots s_1|_{T_w}\text{ mod }K_w)_{w\in N(\xi)}$ can take only finitely many values. Since the latter determines uniquely the endpoint of $\gamma$, and every point of $U$ can be connected to the trivial coset by a path inside $U$, this shows that $|U|<\infty$.  \qedhere
\end{proof}

Recall that given a graph $\Gamma$ and a subset $F$ of vertices, we denote by $\pi_0(\Gamma\setminus F)$ the set of connected components of $\Gamma\setminus F$, and by $\pi_0(\Gamma\setminus F)_\infty$ its subset consisting of infinite connected components. Fix a finite subset $\Sigma\subset\Gamma_{G/H}$ such that $|\pi_0(\Gamma_{G/H}\setminus\Sigma)_\infty|\ge 2$.  For $U\in \pi_0(\Gamma_{G/H}\setminus\Sigma)_\infty$, we set $p'(U):=p(U)\setminus p(\Sigma)$.

\begin{lemma}\label{l-projection-union-cc}

For every $U\in \pi_0(\Gamma_{G/H}\setminus\Sigma)_\infty$, $p'(U)$ is non-empty and it is a union of elements of $\pi_0(\Gammatilde_\xi\setminus p(\Sigma))$.
    
\end{lemma}
\begin{proof}
    The fact that $p'(U)$ is non-empty follows from Lemma \ref{l-fiber-disconnected}, since $p(\Sigma)$ is finite. Let $x$ be any point in $p'(U)$, and $Q\in \pi_0(\Gammatilde_\xi\setminus p(\Sigma))$ be its connected component. For any point $y\in Q$, we can find a path $\gamma\subset Q$ starting at $x$ and ending at $y$. The lift of $\gamma$ starting at any point in $p^{-1}(x)\cap U$ avoids $\Sigma$ (since it projects to $\gamma$ which avoids $p(\Sigma)$), hence it is entirely contained in $U$. Since its endpoint projects to $y$, it follows that $y\in p'(U)$. Since $y\in Q$ is arbitrary this shows that $Q\subset p'(U)$. We have shown that $p'(U)$ contains any component  $Q\in \pi_0(\Gammatilde_\xi\setminus p(\Sigma))$ that it intersects non-trivially, and so it is a union of such components. \qedhere


    
\end{proof}

\begin{lemma}[Second key step] \label{l-ends-injection}
    Let $U_1, U_2\in \pi_0(\Gamma_{G/H}\setminus\Sigma)_\infty$ be such that $p'(U_1)\cap p'(U_2)$ contains some infinite component $Q\in \pi_0(\Gammatilde_\xi\setminus p(\Sigma))_\infty$. Then $U_1=U_2$.
\end{lemma}
\begin{proof}
    Fix $xH\in U_1, yH\in U_2$ such that $p(xH)=p(yH)\in Q$. Then $xG^0_\xi=yG^0_\xi$ and so we can rewrite $y=xk$ where $k=x^{-1}y\in G^0_\xi$. We fix such $x$ and $k$ for the rest of the proof. Since $k\in G^0_\xi$, we can choose a prefix $w$ of $\xi$ such that $k\in G_w^0$. Consider the subgroup
\[R=\langle K_v \colon v\in N(\xi) \colon |v|\le |w|\rangle\le G_w^0.\]
(where the $K_v\leq \rist(v)$ are still the subgroups from \eqref{e-subgroup}). Then $R$ has finite index in $G_w^0$.  Recall that in a finitely generated branch group $G$, all rigid stabilizers $\rist(v)$ are finitely generated, and hence so are the groups $K_v$, hence the group $R$ is finitely generated. As a consequence, the group  $G_w^0$ is also finitely generated. Therefore, there is a finite index subgroup $R_0<R$ which is characteristic in $G_w^0$, and thus normal in the whole vertex stabilizer $G_w$.

Let $\Omega\subset G_w^0$ be a set of representatives of the classes of the finite quotient $G_w^0/R_0$. Set $C=\sup\{\ell_S(t), t\in \Omega\}$, where $\ell_S(\cdot)$ denotes the word length.

We have $G^0_\xi\le G_w$, and the set of cosets $G_w/G^0_\xi$ defines a subset of the set of vertices  of the graph of germs $\Gammatilde_\xi$. Since $G_w$ has finite index in $G$, there exists some $D>0$ such that every vertex of $\Gammatilde_\xi$ is at distance at most $D$ from $G_w/G^0_\xi$. On the other hand, since $Q$ is a connected component of $\Gammatilde_\xi \setminus p(\Sigma)$, the ball of radius $D$ around all but finitely many points of $Q$ is contained in $Q$. In particular, $Q$ has infinite intersection with $G_w/G^0_\xi$.  Therefore we can find a path $\gamma$ in $Q$ that starts at $xG_\xi^0=p(xH)=p(xkH)$ and ends at  a point  of $G_w/G^0_\xi$ at distance $>C$ from $p(\Sigma)$. Let $h=s_n\cdots s_1$ be the word read along this path, so that its endpoint is $zG_\xi^0$, where $z=hx\in G_w$. Then the two lifts of $\gamma$ starting at $xH$ and $xkH$ both avoid $\Sigma$, and thus are contained in $U_1$ and $U_2$, respectively. In particular $zH\in U_1$ and $zkH\in U_2$.  Furthermore both points $zH$ and $zkH$ are at distance $>C$ from $\Sigma$. Let $t\in \Omega$ be a representative of $zk^{-1}z^{-1}$ in $G_w^0/R_0$. Then there is a path in $\Gamma_{G/H}$ of length at most $C$ going from $zkH$ to $tzkH$. This path avoids $\Sigma$, which implies that $tzkH\in U_2$. We have $zkz^{-1}t\in R_0$, by the choice of $t$. Since $k, z\in G_w$, and $R_0$ is normal in $G_w$, this implies that $z^{-1}tzk\in R_0\le H$, so
\[tzkH=z\cdot z^{-1}tzkH=zH\in U_1.\]
We have shown that $tzkH\in U_1\cap U_2\neq \varnothing$,
which implies that $U_1=U_2$. \qedhere
    
\end{proof}

Lemma \ref{l-ends-injection} immediately implies that $|\pi_0(\Gammatilde_\xi\setminus p(\Sigma))_\infty|\ge |\pi_0(\Gamma_{G/H}\setminus \Sigma)_\infty|$,
hence $\esf(\Gammatilde_\xi)\ge\esf(\Gamma_{G/H})\ge 2$. Let us argue that it also implies that $H$ has finite index in $G^0_\xi$ (and hence in $G_\xi$). Suppose by contradiction that $p\colon \Gamma_{G/H}\to \Gammatilde_\xi$ has infinite degree. Let $\gamma$ be a finite path in $\Gammatilde_\xi$ that visits all components in $\pi_0(\Gammatilde_\xi\setminus p(\Sigma))_\infty$. If $p$ has infinite degree, the path $\gamma$ admits infinitely many lifts to $\Gamma_{G/H}$. Hence it admits a lift $\tilde{\gamma}$ which avoids $\Sigma$ as well as all the (finitely many) finite components in $\pi_0(\Gamma_{G/H}\setminus \Sigma)$.  Therefore $\tilde{\gamma}$ is contained in a single infinite component $U\in \pi_0(\Gamma_{G/H}\setminus \Sigma)_\infty$, which then has the property that $p'(U)$ contains all components in $\pi_0(\Gammatilde_\xi\setminus p(\Sigma))_\infty$. Therefore Lemma \ref{l-ends-injection} implies that $U$ is the unique infinite component of $\Gamma_{G/H}\setminus \Sigma$, which contradicts our choice of $\Sigma$.

Finally let $\nu_A$ be an unbounded cardinal-definite function associated to a commensurated subset $A\subset G/H$. Upon changing $A$ in its commensurability class (which changes $\nu_A$ by a bounded amount) we can find a finite subset $\Sigma\subset \Gamma_{G/H}$ such that $A$ is a union of elements of $\pi_0(\Gamma_{G/H}\setminus \Sigma)_
\infty$. Then Lemma \ref{l-ends-injection} implies that  $|A\triangle p^{-1}(p'(A))|<\infty$, thus $\nu_A-\nu_{p^{-1}(p'(A))}$ is bounded. But $\nu_{p^{-1}(p'(A))}=[G^0_\xi \colon H]\nu_{p'(A)}$ by Lemma \ref{lem:endsofcoverings}, finishing the proof. \qedhere

\end{proof}

\end{document}
